\section{Scheduling de Carga en el Depósito}

Antes de iniciar las rutas de distribución, cada camión debe ser cargado en el depósito central. Para este proceso se tiene en consideración que cada compartimento se debe cargar de manera secuencial dentro de cada camión, es decir, el segundo compartimento solo puede comenzar su carga una vez que el primero ha terminado.

Además, el depósito cuenta con una bahía disponible para cada tipo de combustible: una bahía para Regular y una bahía para Diésel. Por lo tanto, no pueden cargarse simultáneamente dos operaciones del mismo producto, pero sí es posible cargar Regular en un camión y Diésel en otro al mismo tiempo.

En resumen, para cargar un camión de un tipo de combustible, debe terminarse la carga anterior para inciar la siguiente y si es que el otro camión no esta cargando su compartimiento del mismo tipo de combustible del que se quiere cargar el primero.

\subsection{Datos de Carga}

Para el escenario determinista se consideran los siguientes volúmenes fijos de carga:

\begin{table}[H]
\centering
\caption{Volúmenes fijos de carga por camión y compartimento}
\begin{tabular}{cccc}
\toprule
Camión & Compartimento & Producto & Volumen (L) \\
\midrule
T1 & C0 & Regular & 5000 \\
T1 & C1 & Diésel & 5000 \\
T2 & C0 & Diésel & 4000 \\
T2 & C1 & Regular & 3500 \\
\bottomrule
\end{tabular}
\end{table}

Las tasas de carga son de 500 L/min para Regular y 400 L/min para Diésel. Además, cada operación de carga considera un tiempo de limpieza de 5 minutos al finalizar.

\subsection{Cálculo de Tiempos de Carga}

El tiempo total de carga de una operación se calcula como:

\begin{equation}
    p_{kc} = \frac{Vol_{kc}}{r_p} + \mu
\end{equation}

Donde:

\begin{itemize}
    \item $p_{kc}$: tiempo total de carga del compartimento $c$ del camión $k$.
    \item $Vol_{kc}$: volumen cargado en el compartimento $c$ del camión $k$.
    \item $r_p$: tasa de carga del producto $p$.
    \item $\mu$: tiempo de limpieza, con $\mu = 5$ minutos.
\end{itemize}

Aplicando la fórmula anterior, se obtienen los siguientes tiempos:

\begin{table}[H]
\centering
\caption{Tiempos de carga por operación}
\begin{tabular}{cccccc}
\toprule
Operación & Producto & Volumen (L) & Tasa (L/min) & Limpieza (min) & Tiempo total (min) \\
\midrule
T1-C0 & Regular & 5000 & 500 & 5 & 15 \\
T1-C1 & Diésel & 5000 & 400 & 5 & 17.5 \\
T2-C0 & Diésel & 4000 & 400 & 5 & 15 \\
T2-C1 & Regular & 3500 & 500 & 5 & 12 \\
\bottomrule
\end{tabular}
\end{table}

\subsection{Conjuntos y Parámetros del Scheduling}

\begin{itemize}
    \item $O$: conjunto de operaciones de carga.
    \item $B$: conjunto de bahías, donde $B = \{R,D\}$.
    \item $p_o$: duración de la operación de carga $o$.
    \item $b_o$: bahía requerida por la operación $o$.
    \item $M$: constante suficientemente grande.
\end{itemize}

Cada operación $o \in O$ representa la carga de un compartimento específico de un camión. Para este caso:

\begin{equation*}
O = \{T1C0, T1C1, T2C0, T2C1\}
\end{equation*}

\subsection{Variables de Decisión del Scheduling}

\begin{itemize}
    \item $s_o \geq 0$: tiempo de inicio de la operación $o$.
    \item $f_o \geq 0$: tiempo de término de la operación $o$.
    \item $C_{\max} \geq 0$: makespan del proceso de carga.
    \item $y_{oo'} \in \{0,1\}$: vale 1 si la operación $o$ se realiza antes que la operación $o'$ en la misma bahía.
\end{itemize}

\subsection{Función Objetivo del Scheduling}

El objetivo del scheduling es minimizar el makespan, es decir, el tiempo en que finaliza la última operación de carga.

\begin{equation}
    \min C_{\max}
\end{equation}

\subsection{Restricciones del Scheduling}

\paragraph{Término de cada operación.}

El tiempo de término de cada operación corresponde a su tiempo de inicio más su duración.

\begin{equation}
    f_o = s_o + p_o,
    \quad \forall o \in O
\end{equation}

\paragraph{Precedencia dentro de cada camión.}

La carga del segundo compartimento de cada camión solo puede comenzar una vez finalizada la carga del primer compartimento.

\begin{align}
    s_{T1C1} &\geq f_{T1C0} \\
    s_{T2C1} &\geq f_{T2C0}
\end{align}

\paragraph{No solapamiento por bahía.}

Dos operaciones que utilizan la misma bahía no pueden ejecutarse simultáneamente.

\begin{align}
    s_{o'} &\geq f_o - M(1-y_{oo'}),
    && \forall o,o' \in O, o \neq o', b_o = b_{o'} \\
    s_o &\geq f_{o'} - My_{oo'},
    && \forall o,o' \in O, o \neq o', b_o = b_{o'}
\end{align}

\paragraph{Definición del makespan.}

El makespan debe ser mayor o igual al término de cada operación.

\begin{equation}
    C_{\max} \geq f_o,
    \quad \forall o \in O
\end{equation}

\subsection{Análisis de solución propuesta}

Una secuencia factible y eficiente es la siguiente:

\begin{table}[H]
\centering
\caption{Secuencia propuesta de carga por bahía}
\begin{tabular}{ccccc}
\toprule
Bahía & Operación & Inicio (min) & Término (min) & Duración (min) \\
\midrule
Regular & T1-C0 & 0 & 15 & 15 \\
Diésel & T2-C0 & 0 & 15 & 15 \\
Regular & T2-C1 & 15 & 27 & 12 \\
Diésel & T1-C1 & 15 & 32.5 & 17.5 \\
\bottomrule
\end{tabular}
\end{table}

Con esta planificación, ambos camiones respetan la precedencia interna de carga. El camión T1 carga primero su compartimento C0 y luego su compartimento C1. Del mismo modo, el camión T2 carga primero C0 y luego C1.

\subsection{Makespan y Verificación de Horarios de Salida}

El makespan de la secuencia propuesta corresponde al término de la última operación:

\begin{equation}
    C_{\max} = 32.5 \text{ minutos}
\end{equation}

Si se inicia el proceso de carga a a las 04:27:30, el último camión queda listo a las 05:00. Por lo tanto, el camión T1 puede cumplir su hora de salida programada a las 05:00. Además, el camión T2 finaliza su carga luego de 27 minutos desde el inicio del proceso, por lo que estaría listo antes de las 05:30.

En consecuencia, la secuencia propuesta permite verificar y confirmar que ambos camiones pueden salir a sus horas programadas, siempre que el proceso de carga inicie con la anticipación suficiente.


\subsection{Diagrama de Gantt Esquemático}

\begin{table}[H]
\centering
\caption{Diagrama de Gantt esquemático del proceso de carga}
\begin{tabular}{c|cccc}
\toprule
Bahía & 0--15 & 15--27 & 27--32.5 & Makespan \\
\midrule
Regular & T1-C0 & T2-C1 & -- & 27 \\
Diésel & T2-C0 & -- & T1-C1 & 32.5 \\
\bottomrule
\end{tabular}
\end{table}